\documentclass[10pt,a4paper]{article}

\usepackage[T1]{fontenc}
\usepackage[utf8]{inputenc}
\usepackage{lmodern}
\usepackage{microtype}
\usepackage{geometry}
\geometry{left=18mm,right=18mm,top=18mm,bottom=20mm,headheight=20pt}
\usepackage{setspace}
\setstretch{1.02}
\usepackage{amsmath,amssymb,mathtools,bm}
\usepackage{array,tabularx,booktabs}
\usepackage{enumitem}
\usepackage{xcolor}
\usepackage[most]{tcolorbox}
\usepackage{needspace}
\usepackage{etoolbox}
\usepackage{fancyhdr}
\usepackage{titlesec}
\usepackage[hidelinks]{hyperref}

\definecolor{ink}{HTML}{414141}
\definecolor{vuBurgundy}{HTML}{78003F}
\definecolor{vuPink}{HTML}{E64164}
\definecolor{blue}{HTML}{174F78}
\definecolor{paleBlue}{HTML}{F1F6F9}
\definecolor{paleBurgundy}{HTML}{FAF2F6}
\definecolor{palePink}{HTML}{FDF3F5}
\definecolor{softGray}{HTML}{F5F5F5}
\definecolor{workGray}{HTML}{ABB3BB}

\color{ink}
\setlength{\parindent}{0pt}
\setlength{\parskip}{0.42em}
\raggedbottom
\clubpenalty=10000
\widowpenalty=10000
\displaywidowpenalty=10000
\setlist[itemize]{leftmargin=1.5em,itemsep=0.18em,topsep=0.20em}
\setlist[enumerate]{leftmargin=3.0em,itemsep=0.20em,topsep=0.20em}

\titleformat{\section}
  {\Large\bfseries\color{vuBurgundy}}{\thesection}{0.65em}{}
\titleformat{\subsection}
  {\large\bfseries\color{blue}}{\thesubsection}{0.65em}{}
\titlespacing*{\section}{0pt}{0.9em}{0.50em}
\titlespacing*{\subsection}{0pt}{0.7em}{0.35em}
\pretocmd{\section}{\Needspace{14\baselineskip}}{}{}
\pretocmd{\subsection}{\Needspace{5\baselineskip}}{}{}

\pagestyle{fancy}
\fancyhf{}
\fancyhead[L]{\small\color{vuBurgundy}\textbf{Algebraic structures}}
\fancyhead[R]{\small\color{ink!65}From binary operations to $K$-algebras}
\fancyfoot[C]{\small\color{vuBurgundy}\thepage}
\renewcommand{\headrulewidth}{0.4pt}
\renewcommand{\headrule}{\hbox to\headwidth{%
  \color{vuBurgundy}\leaders\hrule height \headrulewidth\hfill}%
  \vskip-\headrulewidth}

\hypersetup{
  pdftitle={Algebra over a field K},
  pdfsubject={Axiom systems for groups, rings, fields, vector spaces, bilinear maps, and algebras},
  colorlinks=true,
  linkcolor=vuBurgundy,
  urlcolor=blue
}

\newcommand{\R}{\mathbb{R}}
\newcommand{\C}{\mathbb{C}}
\newcommand{\Z}{\mathbb{Z}}
\newcommand{\Q}{\mathbb{Q}}
\newcommand{\N}{\mathbb{N}}
\newcommand{\End}{\operatorname{End}}

\newtcbtheorem[number within=section]{definition}{Definition}{
  enhanced,colback=paleBlue,colframe=blue,
  colbacktitle=blue,coltitle=white,fonttitle=\bfseries,
  boxrule=0.7pt,arc=1.4mm,
  left=2.5mm,right=2.5mm,top=1.5mm,bottom=1.5mm,
  before skip=5pt,after skip=6pt
}{def}

\newtcbtheorem[use counter from=definition]{proposition}{Proposition}{
  enhanced,colback=paleBurgundy,colframe=vuBurgundy,
  colbacktitle=vuBurgundy,coltitle=white,fonttitle=\bfseries,
  boxrule=0.7pt,arc=1.4mm,
  left=2.5mm,right=2.5mm,top=1.5mm,bottom=1.5mm,
  before skip=5pt,after skip=6pt
}{prop}

\newtcbtheorem[use counter from=definition]{examplebox}{Example}{
  enhanced,colback=softGray,colframe=ink!45,
  colbacktitle=ink!65,coltitle=white,fonttitle=\bfseries,
  boxrule=0.6pt,arc=1.4mm,
  left=2.5mm,right=2.5mm,top=1.5mm,bottom=1.5mm,
  before skip=5pt,after skip=6pt
}{ex}

\newtcolorbox{precision}[1][Remark]{
  enhanced,colback=palePink,colframe=vuPink,
  title={#1},colbacktitle=vuBurgundy,coltitle=white,fonttitle=\bfseries,
  boxrule=0.7pt,arc=1.4mm,
  left=2.5mm,right=2.5mm,top=1.5mm,bottom=1.5mm,
  before skip=5pt,after skip=6pt
}

\newtcolorbox{summarybox}[1]{
  enhanced,colback=softGray,colframe=workGray,
  title={#1},colbacktitle=ink!70,coltitle=white,fonttitle=\bfseries,
  boxrule=0.65pt,arc=1.4mm,
  left=2.5mm,right=2.5mm,top=1.5mm,bottom=1.5mm,
  before skip=5pt,after skip=6pt
}

\newcommand{\keyidea}[1]{%
  \par\smallskip
  \begin{minipage}{\linewidth}
  \noindent\textbf{\color{vuBurgundy}Key idea.} #1
  \end{minipage}
  \par\smallskip
}

\BeforeBeginEnvironment{proposition}{\Needspace{12\baselineskip}}

\begin{document}

\begin{center}
  {\Huge\bfseries\color{vuBurgundy}Algebra over a field $K$\par}
  \vspace{2mm}
  {{\large Groups, rings, fields, vector spaces, and bilinear maps\par}}
\end{center}

An algebra combines a vector-space structure with a compatible multiplication.
We build toward this definition from binary operations, through groups, rings,
fields, and vector spaces. Linear and bilinear maps then provide the language
needed to formulate the compatibility precisely. Important geometric examples
are the algebras of smooth functions and smooth function germs.

The labels attached to axioms identify the structure under discussion. Thus
$(G1)$ denotes a group axiom, $(R1)$ denotes a ring axiom, and $(V1)$ denotes a
vector-space axiom. These labels are local conventions rather than universal
names, but they make dependencies between definitions visible.
We write $\N=\{1,2,3,\ldots\}$ and $\N_0=\{0,1,2,\ldots\}$.

\section{Binary operations}

\begin{definition}{Binary operation}{binary-operation}
Let $S$ be a set. A \emph{binary operation on $S$} is a function
\[
  \ast:S\times S\longrightarrow S,
  \qquad
  (a,b)\longmapsto a\ast b.
\]
The codomain $S$ means that $S$ is closed under the operation $\ast$.
In other words, if $a,b\in S$, then $a\ast b\in S$.
\end{definition}

The symbol $\ast$ carries no fixed meaning. It can denote addition,
multiplication, composition, or another operation. What matters is the map and
the properties it satisfies.

\begin{definition}{Standard properties of a binary operation}{operation-properties}
Let $\ast$ be a binary operation on $S$. The following properties will be
used repeatedly.
\begin{enumerate}[label=\textup{(O\arabic*)}]
  \item \textbf{Associativity.} For every $a,b,c\in S$,
  \[
    (a\ast b)\ast c=a\ast(b\ast c).
  \]

  \item \textbf{Identity.} There exists $e\in S$ such that, for every
  $a\in S$,
  \[
    e\ast a=a\ast e=a.
  \]

  \item \textbf{Existence of inverses.} Relative to an identity $e$, every
  $a\in S$ has a two-sided inverse, that is, an element $b\in S$ such that
  \[
    a\ast b=b\ast a=e.
  \]

  \item \textbf{Commutativity.} For every $a,b\in S$,
  \[
    a\ast b=b\ast a.
  \]
\end{enumerate}
\end{definition}

An identity, when it exists, is unique. If the operation
is also associative, inverses are unique whenever they exist. In this
setting the notation $a^{-1}$ is unambiguous. These assertions are proved in Proposition~\ref{prop:uniqueness}.

\emph{Distributivity} relates two operations. If $+$ and $\cdot$ are binary
operations on the same set, then multiplication distributes over addition on
the left and on the right, respectively, when
\[
  a(b+c)=ab+ac,
  \qquad
  (a+b)c=ac+bc
\]
for all $a,b,c$. Both laws will be required in the definition of a ring.

\begin{precision}[Closure is built into the definition]
One sometimes sees closure listed as a separate axiom. In the notation
$\ast:S\times S\to S$, closure is already encoded by the codomain. A rule
such as subtraction on $\N$ is not a binary operation on $\N$ because the
result can lie outside $\N$.
\end{precision}

\begin{examplebox}{Composition of maps}{composition}
Let $X$ be a set and let $S$ be the set of all functions $X\to X$. Composition
defines a binary operation on $S$ by
\[
  (f,g)\longmapsto f\circ g.
\]
It is associative. It has the identity map $\operatorname{id}_X$ as an
identity element. It is generally not commutative.
\end{examplebox}

\keyidea{A binary operation is a function with two inputs from the same set
and one output in that set. Its additional properties are separate statements
that must be proved or assumed.}

\section{Semigroups, monoids, and groups}

The next structures are obtained by imposing progressively stronger axioms on
one binary operation.

\begin{definition}{Semigroup}{semigroup}
A \emph{semigroup} is a nonempty set $S$ with a binary operation $\ast$ satisfying
the following axiom.
\begin{enumerate}[label=\textup{(S\arabic*)}]
  \item \textbf{Associativity.} For every $a,b,c\in S$,
  \[
    (a\ast b)\ast c=a\ast(b\ast c).
  \]
\end{enumerate}
\end{definition}

\begin{definition}{Monoid}{monoid}
A \emph{monoid} is a set $M$ with a binary operation $\ast$ satisfying
the following axioms.
\begin{enumerate}[label=\textup{(M\arabic*)}]
  \item \textbf{Associativity.} For every $a,b,c\in M$,
  \[
    (a\ast b)\ast c=a\ast(b\ast c).
  \]

  \item \textbf{Identity.} There exists $e\in M$ such that, for every
  $a\in M$,
  \[
    e\ast a=a\ast e=a.
  \]
\end{enumerate}
\end{definition}

\begin{definition}{Group and abelian group}{group}
A \emph{group} is a set $G$ with a binary operation $\ast$ satisfying
the following axioms.
\begin{enumerate}[label=\textup{(G\arabic*)}]
  \item \textbf{Associativity.} For every $a,b,c\in G$,
  \[
    (a\ast b)\ast c=a\ast(b\ast c).
  \]

  \item \textbf{Identity.} There exists $e\in G$ such that, for every
  $a\in G$,
  \[
    e\ast a=a\ast e=a.
  \]

  \item \textbf{Inverses.} For every $a\in G$, there exists $a^{-1}\in G$
  such that
  \[
    a\ast a^{-1}=a^{-1}\ast a=e.
  \]
\end{enumerate}
The group is \emph{abelian} if it also satisfies
\begin{enumerate}[label=\textup{(G4)},leftmargin=3em]
  \item \textbf{Commutativity.} For every $a,b\in G$,
  \[
    a\ast b=b\ast a.
  \]
\end{enumerate}
\end{definition}

The definitions give the implication chain
\[
  \text{abelian group}
  \Longrightarrow
  \text{group}
  \Longrightarrow
  \text{monoid}
  \Longrightarrow
  \text{semigroup}.
\]
Each arrow means that the structure on the left satisfies every axiom required
by the structure on the right. The reverse implications fail in general.

\begin{proposition}{Uniqueness of the identity and inverses}{uniqueness}
A binary operation has at most one two-sided identity. If the operation is
associative and has an identity, each element has at most one two-sided
inverse. In particular, a group's identity and inverses are unique.
\end{proposition}

\textit{Proof.}
Suppose $e$ and $e'$ are both identities. Then
\[
  e=e\ast e'=e'.
\]
If the operation is associative and $b$ and $c$ are both inverses of $a$, then
\[
  b=b\ast e
   =b\ast(a\ast c)
   =(b\ast a)\ast c
   =e\ast c=c.
\]
Thus identities are unique, and associativity ensures uniqueness of inverses
whenever they exist.
\hfill$\square$

\begin{examplebox}{Additive examples}{additive-groups}
The pair $(\N_0,+)$ is a commutative monoid but not a group. A positive integer
has no additive inverse in $\N_0$.

The pair $(\Z,+)$ is an abelian group. Its identity is $0$, and the inverse of
$n\in\Z$ is $-n$.
\end{examplebox}

\section{Rings}

A ring places two binary operations on one set. Their usual names are addition
and multiplication. The notation alone does not imply any axioms.

\begin{definition}{Ring}{ring}
A \emph{ring} is a set $A$ with two binary operations
\[
  +:A\times A\longrightarrow A,
  \qquad
  \cdot:A\times A\longrightarrow A,
\]
satisfying the following axioms.
\begin{enumerate}[label=\textup{(R\arabic*)}]
  \item \textbf{Additive abelian group.} The additive structure $(A,+)$ is
  an abelian group.

  \item \textbf{Associative multiplication.} For every $a,b,c\in A$,
  \[
    (ab)c=a(bc).
  \]

  \item \textbf{Distributivity.} For every $a,b,c\in A$,
  \[
    (a+b)c=ac+bc,
    \qquad
    a(b+c)=ab+ac.
  \]
\end{enumerate}
\end{definition}

Axiom $(R1)$ includes all four group axioms for addition. Thus addition is
associative and commutative, $A$ has an additive identity $0_A$, and every
$a\in A$ has an additive inverse $-a$. We write $ab$ for $a\cdot b$.

Since $A$ is nonempty by $(R1)$, axiom $(R2)$ makes $(A,\cdot)$ a
semigroup. Axiom $(R3)$ relates multiplication to addition. Without commutativity
of multiplication, neither distributive law follows from the other, so both are
stated.

\begin{definition}{Additional ring properties}{ring-properties}
A ring may satisfy either of the following additional axioms.
\begin{enumerate}[label=\textup{(R\arabic*)},start=4]
  \item \textbf{Multiplicative identity.} There exists $1_A\in A$ such
  that, for every $a\in A$,
  \[
    1_Aa=a1_A=a.
  \]
  A ring satisfying $(R4)$ is called \emph{unital}.

  \item \textbf{Commutative multiplication.} For every $a,b\in A$,
  \[
    ab=ba.
  \]
  A ring satisfying $(R5)$ is called \emph{commutative}.
\end{enumerate}
\end{definition}

If $(R4)$ holds, then $(A,\cdot)$ is a monoid. The zero element and the unit
must not be confused. The \emph{zero ring} consists of the single element
$0_A$. Its multiplicative identity is also $0_A$. In every other unital ring,
$0_A\neq1_A$. Indeed, $1_A=0_A$ would imply $a=1_Aa=0_Aa=0_A$ for every $a$, by
the following proposition.

\begin{proposition}{Elementary consequences of the ring axioms}{ring-zero}
For every $a,b\in A$,
\[
  0_Aa=a0_A=0_A,
  \qquad
  (-a)b=-(ab),
  \qquad
  a(-b)=-(ab).
\]
\end{proposition}

\textit{Proof.}
Using distributivity,
\[
  0_Aa=(0_A+0_A)a=0_Aa+0_Aa.
\]
Adding the additive inverse of $0_Aa$ to both sides gives $0_Aa=0_A$.
The proof of $a0_A=0_A$ is analogous. Also,
\[
  ab+(-a)b=(a+(-a))b=0_Ab=0_A,
\]
so $(-a)b$ is the additive inverse of $ab$. The last identity is analogous.
\hfill$\square$

\begin{examplebox}{Integers and matrices}{ring-examples}
The integers $\Z$ form a commutative unital ring.

For a field $K$ and $n\geq1$, let $M_n(K)$ denote the set of $n\times n$
matrices with entries in $K$. It is a unital ring under matrix addition and
matrix multiplication, with unit $I_n$. For $n\geq2$, it is noncommutative.
For example, in $M_2(K)$,
\[
  P=\begin{pmatrix}1&0\\0&0\end{pmatrix},
  \qquad
  Q=\begin{pmatrix}0&1\\0&0\end{pmatrix},
  \qquad PQ=Q,\quad QP=0.
\]
The same example embeds into larger matrices by filling the other entries
with zeros.
\end{examplebox}

\medskip
\noindent\textbf{\color{vuBurgundy}Convention concerning units.}
Here a ring need not have a multiplicative identity. The adjective
\emph{unital} states this additional assumption. Another common convention
includes a multiplicative identity in the definition of a ring, so the chosen
convention should always be kept in mind.


\section{Fields}

A field is a ring in which the nonzero elements also form an abelian group
under multiplication.

\begin{definition}{Field}{field}
A \emph{field} is a set $K$ with binary operations
$+,\cdot:K\times K\to K$ satisfying the following axioms.
\begin{enumerate}[label=\textup{(F\arabic*)}]
  \item \textbf{Additive abelian group.} The additive structure $(K,+)$ is
  an abelian group.

  \item \textbf{Multiplicative abelian group.} The multiplicative structure
  $(K\setminus\{0_K\},\cdot)$ is an abelian group.

  \item \textbf{Distributivity.} For every $r,s,t\in K$,
  \[
    r(s+t)=rs+rt,
    \qquad
    (r+s)t=rt+st.
  \]
\end{enumerate}
\end{definition}

Axiom $(F2)$ includes multiplicative associativity, commutativity, the identity
$1_K$, and an inverse $r^{-1}$ for every nonzero $r\in K$. It also implies
$0_K\neq1_K$ because the identity $1_K$ belongs to $K\setminus\{0_K\}$.
Closure of this group asserts that a product of nonzero elements is nonzero.

\begin{proposition}{Equivalent ring formulation of a field}{field-ring}
A field is equivalently a commutative unital ring $K$ with $1_K\neq0_K$ in
which every nonzero element has a multiplicative inverse.
\end{proposition}

\textit{Proof.}
Assume $(F1)$--$(F3)$. Additive cancellation and distributivity give
$0_Kr=r0_K=0_K$ by the same argument as in
Proposition~\ref{prop:ring-zero}. That argument does not use multiplicative
associativity. By $(F2)$, multiplication is associative and commutative on
the nonzero elements and has identity $1_K$. The identities
$0_Kr=r0_K=0_K$ then show that associativity, commutativity, and the identity
law also hold whenever one of the factors is $0_K$. Hence these properties hold
on all of $K$. Thus $K$ is a commutative unital ring with $1_K\neq0_K$, and
$(F2)$ supplies the inverses.

Conversely, let $K$ be such a ring. If $r,s\neq0_K$ and $rs=0_K$, then
\[
  s=(r^{-1}r)s=r^{-1}(rs)=0_K,
\]
a contradiction. Therefore the nonzero elements are closed under
multiplication. An inverse of a nonzero element is itself nonzero, since
$r0_K=0_K\neq1_K$. The ring laws and the assumed inverses now make
$K\setminus\{0_K\}$ an abelian group. The other field axioms are ring axioms.
\hfill$\square$

The condition $1_K\neq0_K$ is essential. In the zero ring there are no nonzero
elements, so the requirement that every nonzero element have an inverse would
hold vacuously.

\begin{examplebox}{Fields and nonfields}{field-examples}
The sets
\[
  \Q,
  \qquad
  \R,
  \qquad
  \C
\]
are fields with their usual operations.

The integers $\Z$ are not a field because $2^{-1}=1/2$ does not belong to
$\Z$.

For every prime number $p$, the residue classes $\mathbb F_p=\Z/p\Z$ form a
finite field.
\end{examplebox}

\begin{summarybox}{Where the group structures occur}
\renewcommand{\arraystretch}{1.25}
\begin{tabularx}{\linewidth}{>{\bfseries}p{3.2cm}X X}
\toprule
Structure & Additive part & Multiplicative part\\
\midrule
Ring $A$ & $(A,+)$ is an abelian group & $(A,\cdot)$ is a semigroup\\
Unital ring $A$ & $(A,+)$ is an abelian group & $(A,\cdot)$ is a monoid\\
Field $K$ & $(K,+)$ is an abelian group & $(K\setminus\{0\},\cdot)$ is an abelian group
\\\bottomrule
\end{tabularx}
\end{summarybox}

\keyidea{The field $K$ used in linear algebra is itself an algebraic
structure. Its elements become the scalars acting on a vector space.}

\section{Vector spaces over a field}

Fix a field $K$. A vector-space structure consists of an addition on a set
$V$ and an action of the scalars in $K$ on the elements of $V$.

\begin{definition}{$K$-vector space}{vector-space}
A \emph{vector space over $K$} is a set $V$ equipped with vector addition
\[
  +:V\times V\longrightarrow V,
  \qquad
  (u,v)\longmapsto u+v,
\]
and scalar multiplication
\[
  K\times V\longrightarrow V,
  \qquad
  (r,v)\longmapsto rv,
\]
satisfying the following axioms for all $r,s\in K$ and $u,v\in V$.
\begin{enumerate}[label=\textup{(V\arabic*)}]
  \item \textbf{Additive abelian group.} The additive structure $(V,+)$ is
  an abelian group.

  \item \textbf{Distributivity over vector addition.}
  \[
    r(u+v)=ru+rv.
  \]

  \item \textbf{Distributivity over scalar addition.}
  \[
    (r+s)v=rv+sv.
  \]

  \item \textbf{Compatibility with multiplication in $K$.}
  \[
    (rs)v=r(sv).
  \]

  \item \textbf{Scalar identity.}
  \[
    1_Kv=v.
  \]
\end{enumerate}
\end{definition}

The phrase \emph{over $K$} identifies the field of scalars. It does not mean
that $V$ is a subset of $K$.

\begin{summarybox}{The two operations of a vector space}
\renewcommand{\arraystretch}{1.25}
\begin{tabularx}{\linewidth}{>{\bfseries}p{4.2cm}p{4.2cm}X}
Vector addition & $V\times V\to V$ & $(u,v)\mapsto u+v$\\
Scalar multiplication & $K\times V\to V$ & $(r,v)\mapsto rv$
\end{tabularx}
\end{summarybox}

Scalar multiplication is called an \emph{external operation}. It takes a
scalar from $K$ and a vector from $V$. In general, it is not a binary operation
on $V$. The roles of $K$ and $V$ remain distinct even when the underlying sets
coincide, as in the vector space $K$ over itself. Similarly, $+$ denotes both
scalar addition and vector addition. Its meaning is determined by the inputs.

\begin{examplebox}{Standard vector spaces}{vector-examples}
The Cartesian space $K^n$ is a vector space over $K$ with componentwise
addition and scalar multiplication.

The polynomial space $K[x]$ is a vector space over $K$ with the usual
coefficientwise operations. Its vectors are polynomials of arbitrary finite
degree. Constant polynomials identify a copy of $K$ inside this space.

The matrix space $M_n(\C)$ is a vector space over both $\C$ and $\R$. Its
dimensions are
\[
  \dim_{\C}M_n(\C)=n^2,
  \qquad
  \dim_{\R}M_n(\C)=2n^2.
\]
\end{examplebox}

\begin{proposition}{Elementary consequences of the vector-space axioms}{vector-zero}
For every $r\in K$ and $v\in V$,
\[
  0_Kv=0_V,
  \qquad
  r0_V=0_V,
  \qquad
  (-1_K)v=-v.
\]
\end{proposition}

\textit{Proof.}
The first identity follows from
\[
  0_Kv=(0_K+0_K)v=0_Kv+0_Kv.
\]
The second follows similarly from $r(0_V+0_V)=r0_V+r0_V$. Finally,
\[
  v+(-1_K)v=1_Kv+(-1_K)v=(1_K-1_K)v=0_Kv=0_V.
\]
Hence $(-1_K)v$ is the additive inverse of $v$.
\hfill$\square$

\section{Linear and bilinear maps}

The multiplication of an algebra must be compatible with the vector-space
structure. This compatibility is expressed through linearity and bilinearity.

\begin{precision}[A vector space need not have an internal product]
The vector-space axioms do not define a multiplication $V\times V\to V$. A
quantum state space $\mathcal H$ is a complex vector space, but two state
vectors do not generally have a product that lies again in $\mathcal H$.
\end{precision}

\begin{definition}{Linear map}{linear-map}
Let $V$ and $W$ be vector spaces over the same field $K$. A map
$T:V\to W$ is \emph{$K$-linear} if it satisfies the following axioms.
\begin{enumerate}[label=\textup{(L\arabic*)}]
  \item \textbf{Additivity.} For every $u,v\in V$,
  \[
    T(u+v)=T(u)+T(v).
  \]

  \item \textbf{Homogeneity.} For every $r\in K$ and $v\in V$,
  \[
    T(rv)=rT(v).
  \]
\end{enumerate}
\end{definition}

The two conditions can be combined into
\[
  T(ru+sv)=rT(u)+sT(v)
\]
for all $r,s\in K$ and $u,v\in V$.

\begin{definition}{Bilinear map}{bilinear-map}
Let $U$, $V$, and $W$ be vector spaces over $K$. A map
\[
  B:U\times V\longrightarrow W
\]
is \emph{$K$-bilinear} if it satisfies the following conditions for all
$r,s\in K$, $u,u_1,u_2\in U$, and $v,v_1,v_2\in V$.
\begin{enumerate}[label=\textup{(BL\arabic*)},leftmargin=3.5em]
  \item \textbf{Linearity in the first argument.} For every fixed $v\in V$,
  the map $u\mapsto B(u,v)$ is $K$-linear. Equivalently,
  \[
    B(ru_1+su_2,v)=rB(u_1,v)+sB(u_2,v).
  \]

  \item \textbf{Linearity in the second argument.} For every fixed $u\in U$,
  the map $v\mapsto B(u,v)$ is $K$-linear. Equivalently,
  \[
    B(u,rv_1+sv_2)=rB(u,v_1)+sB(u,v_2).
  \]
\end{enumerate}
\end{definition}

The word \emph{separately} is important. Bilinearity means linearity in one
argument while the other argument is held fixed. It does not mean that $B$ is
linear as a map from the product vector space $U\times V$ to $W$.
Indeed, bilinearity gives $B(u,0)=B(0,v)=0$. If $B$ were also linear on the
product, whose operations are componentwise, then
\[
  B(u,v)=B\bigl((u,0)+(0,v)\bigr)=B(u,0)+B(0,v)=0.
\]
Thus only the zero bilinear map is also linear on the product vector space.

\begin{examplebox}{Dot product and matrix multiplication}{bilinear-examples}
For $n\geq2$, the Euclidean dot product
\[
  \R^n\times\R^n\longrightarrow\R,
  \qquad
  (u,v)\longmapsto u\cdot v
\]
is bilinear. It is not an internal multiplication on $\R^n$ because its output
lies in $\R$ rather than $\R^n$.

Matrix multiplication
\[
  M_n(K)\times M_n(K)\longrightarrow M_n(K)
\]
is bilinear and internal. It is also associative.
\end{examplebox}

\begin{precision}[The cross product]
The cross product $\R^3\times\R^3\to\R^3$ is bilinear and internal, but it is
not associative. For standard basis vectors $e_1,e_2,e_3$,
\[
  (e_1\times e_1)\times e_2=0,
  \qquad
  e_1\times(e_1\times e_2)=e_1\times e_3=-e_2.
\]
It therefore does not define an associative algebra. Instead, the cross
product defines a \emph{Lie algebra bracket} on $\R^3$: it is bilinear,
antisymmetric, and satisfies the Jacobi identity
\[
  u\times(v\times w)+v\times(w\times u)+w\times(u\times v)=0.
\]
\end{precision}

\begin{precision}[Complex inner products]
The standard complex inner product, with the convention
$\langle u,v\rangle=\sum_j\overline{u_j}v_j$, satisfies
\[
  \langle ru,sv\rangle=\overline r\,s\langle u,v\rangle.
\]
It is conjugate-linear in the first argument and linear in the second, hence
\emph{sesquilinear}, rather than $\C$-bilinear. A complex inner product should
therefore not be substituted for a bilinear multiplication.
\end{precision}

\keyidea{For an algebra, the candidate multiplication must first have the
correct type $A\times A\to A$. It must then be bilinear and associative.}

\section{Algebras over a field}

Let $A$ be a vector space over a field $K$. An algebra structure adds an
internal multiplication to $A$ and requires that multiplication to respect the
linear structure. Thus an algebra combines the additive and scalar structure of
a vector space with an internal product.

In this chapter, \emph{algebra} means \emph{associative algebra} unless stated
otherwise. A multiplicative identity is not required, so the word
\emph{algebra} alone does not imply unitality. More generally, one may study
vector spaces with bilinear products that are not associative, as the cross
product illustrates.

\begin{definition}{Associative algebra over $K$}{algebra}
An \emph{associative algebra over $K$} is a $K$-vector space $A$ equipped with
a multiplication map
\[
  \mu:A\times A\longrightarrow A,
  \qquad
  \mu(a,b)=ab,
\]
satisfying the following axioms.
\begin{enumerate}[label=\textup{(A\arabic*)}]
  \item \textbf{Associativity.} For every $a,b,c\in A$,
  \[
    (ab)c=a(bc).
  \]

  \item \textbf{Distributivity (additivity in both arguments).}
  For every $a,b,c\in A$,
  \[
    (a+b)c=ac+bc,
    \qquad
    a(b+c)=ab+ac.
  \]

  \item \textbf{Homogeneity in both arguments.} For every $r\in K$ and
  $a,b\in A$,
  \[
    (ra)b=r(ab),
    \qquad
    a(rb)=r(ab).
  \]
\end{enumerate}
\end{definition}

\keyidea{A $K$-algebra combines a $K$-vector space with an associative internal
$K$-bilinear multiplication $A\times A\to A$.}

Axioms $(A2)$ and $(A3)$ together say that the multiplication map $\mu$ is
$K$-bilinear. Thus the definition can be compressed to one sentence.

\begin{summarybox}{Bilinear formulation}
An associative $K$-algebra is a $K$-vector space equipped with an associative
$K$-bilinear multiplication $A\times A\to A$.
\end{summarybox}

In particular, homogeneity in the two arguments implies
\[
  (ra)(sb)=(rs)(ab)
  \qquad (r,s\in K,\ a,b\in A).
\]
This does not impose commutativity of the algebra product. The order of $a$
and $b$ is unchanged.

\begin{definition}{Additional algebra properties}{algebra-properties}
An associative $K$-algebra may satisfy either of the following additional
axioms.
\begin{enumerate}[label=\textup{(A\arabic*)},start=4]
  \item \textbf{Multiplicative identity.} There exists $1_A\in A$ such that
  $1_Aa=a1_A=a$ for every $a\in A$. The algebra is then called \emph{unital}.

  \item \textbf{Commutativity.} One has $ab=ba$ for all $a,b\in A$. The
  algebra is then called \emph{commutative}.
\end{enumerate}
\end{definition}

\begin{proposition}{Equivalent ring formulation}{algebra-ring}
An associative $K$-algebra is equivalently a ring $A$ that is also a
$K$-vector space, using the same addition, and whose ring multiplication
satisfies
\[
  r(ab)=(ra)b=a(rb)
\]
for every $r\in K$ and $a,b\in A$.
\end{proposition}

\textit{Proof.}
Suppose first that $A$ is a $K$-vector space satisfying $(A1)$, $(A2)$, and
$(A3)$. Axiom $(V1)$ says that $(A,+)$ is an abelian group. Axiom $(A1)$ gives
associativity of multiplication. Axiom $(A2)$ gives both distributive laws.
Therefore $A$ is a ring. Axiom $(A3)$ gives the required compatibility with
scalars.

Conversely, suppose that $A$ is a ring and a $K$-vector space with the stated
compatibility. Ring distributivity gives $(A2)$. The compatibility condition
gives $(A3)$. Ring associativity gives $(A1)$. Hence $A$ is an associative
$K$-algebra.
\hfill$\square$

\begin{precision}[How the algebra axioms fit together]
The vector-space structure already makes $(A,+)$ an abelian group.
Associativity $(A1)$ and distributivity $(A2)$ then give the ring structure.
Homogeneity $(A3)$ is the additional condition that ties the ring
multiplication to scalar multiplication by $K$. Thus a $K$-algebra may be
viewed as a ring and a $K$-vector space whose two structures are compatible.
\end{precision}

\section{Examples and nonexamples}

\begin{summarybox}{The three operations on a $K$-algebra}
\renewcommand{\arraystretch}{1.25}
\begin{tabularx}{\linewidth}{>{\bfseries}p{4.1cm}p{4.2cm}X}
Addition in $A$ & $A\times A\to A$ & $(a,b)\mapsto a+b$\\
Multiplication in $A$ & $A\times A\to A$ & $(a,b)\mapsto ab$\\
Scalar multiplication & $K\times A\to A$ & $(r,a)\mapsto ra$
\end{tabularx}
\end{summarybox}

\subsection{Canonical examples}

\begin{examplebox}{The field as an algebra over itself}{field-algebra}
Every field $K$ is a one-dimensional commutative unital algebra over itself.
The algebra multiplication is the field multiplication, and
\[
  \dim_K K=1.
\]
\end{examplebox}

\begin{examplebox}{Polynomial and matrix algebras}{polynomial-matrix}
The polynomial space $K[x]$ is a commutative unital $K$-algebra under ordinary
polynomial multiplication.

The matrix space $M_n(K)$ is a unital $K$-algebra under matrix multiplication.
For $n\geq2$, it is noncommutative.
\end{examplebox}

\begin{examplebox}{An algebra without an identity}{nonunital-algebra}
The polynomials with zero constant term form the $K$-algebra
\[
  xK[x]=\{xf(x):f(x)\in K[x]\}.
\]
This set is closed under addition, scalar multiplication, and polynomial
multiplication. It has no multiplicative identity. Indeed, a nonzero
$e(x)\in xK[x]$ has degree at least one, so $e(x)x$ has degree at least two
and cannot equal $x$. The zero polynomial cannot be an identity either.
\end{examplebox}

\begin{examplebox}{Endomorphism algebra}{endomorphism}
Let $V$ be a vector space over $K$. The space
\[
  \End_K(V)=\{T:V\to V\mid T\text{ is $K$-linear}\}
\]
is a vector space with pointwise operations
\[
  (S+T)(v)=S(v)+T(v),\qquad (rT)(v)=rT(v).
\]
It becomes a unital $K$-algebra under composition,
\[
  (S,T)\longmapsto S\circ T.
\]
Linearity of the operators makes composition bilinear, and composition is
associative. Its identity is $\operatorname{id}_V$. If $\dim_K V=n<\infty$, a
choice of basis gives a $K$-algebra isomorphism
\[
  \End_K(V)\cong M_n(K).
\]
This identification depends on the chosen basis.
\end{examplebox}

\subsection{Useful nonexamples}

\begin{itemize}
  \item A vector space without a specified internal multiplication is not yet
  an algebra.

  \item For $n\geq2$, the dot product on $\R^n$ is not an algebra multiplication because it
  takes values in $\R$ rather than $\R^n$.
  For $n=1$, under the usual identification $\R^1=\R$, it is ordinary
  multiplication and does define an algebra.

  \item The cross product on $\R^3$ is bilinear and internal, but it is not
  associative.

  \item For $n\geq2$, Hermitian $n\times n$ matrices form a real vector space
  but are not closed under ordinary matrix multiplication. Here
  $H^\dagger=\overline H^{\mathsf T}$, and $H$ is Hermitian when $H^\dagger=H$.
  For Hermitian $H,J$, one has $(HJ)^\dagger=JH$, so $HJ$ is Hermitian exactly
  when $H$ and $J$ commute. For example,
  \[
    H=\begin{pmatrix}1&0\\0&-1\end{pmatrix},\qquad
    J=\begin{pmatrix}0&1\\1&0\end{pmatrix},\qquad
    HJ=\begin{pmatrix}0&1\\-1&0\end{pmatrix}.
  \]
  The first two matrices are Hermitian, whereas their product is not.
\end{itemize}

\section{Smooth functions, germs, and algebra homomorphisms}

The algebraic structures above arise naturally in differential geometry.
Smooth functions provide a basic example, and passing to germs records only
their local behaviour near a point.

\begin{examplebox}{The algebra of smooth functions}{smooth-functions}
For a nonempty open set $U\subseteq\R^n$, let $C^\infty(U)$ denote the set of
smooth real-valued functions on $U$. Define
\[
  (f+g)(x)=f(x)+g(x),\qquad
  (fg)(x)=f(x)g(x),\qquad
  (rf)(x)=rf(x),
\]
for $x\in U$ and $r\in\R$. These functions are smooth, and the algebra
identities follow by evaluating them at each $x$ and using the corresponding
identities in $\R$. Thus $C^\infty(U)$ is a commutative unital $\R$-algebra.
Its zero and identity are the constant functions $0$ and $1$.
\end{examplebox}

\begin{definition}{Smooth function germ}{smooth-germ}
Fix $p\in\R^n$. Consider pairs $(U,f)$, where $U$ is an open neighbourhood
of $p$ and $f\in C^\infty(U)$. Set
\[
  (U,f)\sim(V,g)
  \quad\Longleftrightarrow\quad
  f|_W=g|_W\ \text{on some open neighbourhood }W\subseteq U\cap V\text{ of }p.
\]
This is an equivalence relation. Its equivalence classes are the
\emph{smooth function germs at $p$}. The class of $(U,f)$ is denoted by
$[f]_p$. The set of these classes is denoted by $C_p^\infty(\R^n)$.
\end{definition}

Reflexivity and symmetry follow directly from equality of functions.
For transitivity, if $f=g$ on a neighbourhood $W_1$ and $g=h$ on a
neighbourhood $W_2$, then $f=h$ on $W_1\cap W_2$, which is again an open
neighbourhood of $p$. A germ records agreement near $p$, not merely the
value at $p$. For example, $f(x)=0$ and $g(x)=x$ on $\R$ have the same value
at $0$ but distinct germs there.

\begin{proposition}{The algebra of smooth germs}{germ-algebra}
The set $C_p^\infty(\R^n)$ is a commutative unital $\R$-algebra with
\[
  [f]_p+[g]_p=[f+g]_p,\qquad
  [f]_p[g]_p=[fg]_p,\qquad
  r[f]_p=[rf]_p.
\]
When $f$ and $g$ have different domains, the sum and product are formed on
their intersection. The zero and identity are $[0]_p$ and $[1]_p$.
\end{proposition}

\textit{Proof.}
Suppose $[f]_p=[\widetilde f]_p$ and $[g]_p=[\widetilde g]_p$. Choose open
neighbourhoods $W_f$ and $W_g$ on which the respective representatives agree.
On $W_f\cap W_g$,
\[
  f+g=\widetilde f+\widetilde g,\qquad
  fg=\widetilde f\,\widetilde g,\qquad
  rf=r\widetilde f.
\]
Consequently the resulting germs do not depend on the representatives.
Every vector-space and algebra identity now follows from the corresponding
identity for smooth functions, after restricting the finitely many
representatives to a common neighbourhood.\hfill$\square$

\begin{definition}{Algebra homomorphism}{algebra-homomorphism}
Let $A$ and $B$ be $K$-algebras. A \emph{$K$-algebra homomorphism} is a
$K$-linear map $\varphi:A\to B$ that preserves multiplication,
\[
  \varphi(ab)=\varphi(a)\varphi(b)\qquad(a,b\in A).
\]
If $A$ and $B$ are unital, the homomorphism is called \emph{unital} when it
also satisfies $\varphi(1_A)=1_B$.
\end{definition}

For unital algebras, preservation of the multiplicative identity is an
additional condition in the convention used here. It does not follow from
linearity and preservation of products. For example, the zero map
$\R\to\R$ has both properties but does not preserve $1$.

\begin{examplebox}{Evaluation at a point}{evaluation}
The map
\[
  \operatorname{ev}_p:C_p^\infty(\R^n)\longrightarrow\R,
  \qquad [f]_p\longmapsto f(p),
\]
is well defined because equivalent representatives agree at $p$. It is
$\R$-linear, satisfies
$\operatorname{ev}_p([f]_p[g]_p)=f(p)g(p)$, and sends $[1]_p$ to $1$.
Hence it is a unital algebra homomorphism.
\end{examplebox}

\begin{precision}[Connection with derivations at a point]
The product in the germ algebra is the product that appears in the Leibniz
rule. A \emph{derivation at $p$} is an $\R$-linear map
$D:C_p^\infty(\R^n)\to\R$ satisfying
\[
  D([f]_p[g]_p)=D([f]_p)\,g(p)+f(p)\,D([g]_p).
\]
For example, $D([f]_p)=\sum_{i=1}^n v^i\,\partial_i f(p)$ is a derivation
for each $v\in\R^n$: it is well defined on germs, and the ordinary product
rule gives the displayed identity. Evaluation preserves products, whereas
a derivation differentiates them by this rule.
\end{precision}

\end{document}
